\textbf{Phần 1. Trắc nghiệm lựa chọn}
\setcounterpageref{ex}{0}

\begin{ex} %[5H5V25]
Một hình tam giác có diện tích $45\text{ cm}^2$ và chiều cao bằng $7{,}5\text{ cm}$. Độ dài cạnh đáy tương ứng của hình tam giác đó là
\choice
{$6\text{ cm}$}
{\True $12\text{ cm}$}
{$15\text{ cm}$}
{$9\text{ cm}$}
\loigiai{
Độ dài cạnh đáy của hình tam giác là:
$45 \times 2 : 7{,}5 = 90 : 7{,}5 = 12\text{ (cm)}$.\\
\textbf{Đáp án của câu hỏi là B}
}
\end{ex}

\begin{ex} %[5H5V25]
Một tam giác vuông có hai cạnh góc vuông là $12\text{ cm}$ và $16\text{ cm}$, cạnh huyền dài $20\text{ cm}$. Chiều cao tương ứng với cạnh huyền là
\choice
{$8\text{ cm}$}
{$9\text{ cm}$}
{\True $9{,}6\text{ cm}$}
{$10{,}2\text{ cm}$}
\loigiai{
Diện tích của tam giác vuông là:
$12 \times 16 : 2 = 96\text{ (cm}^2\text{)}$.\\
Chiều cao tương ứng với cạnh huyền là:
$96 \times 2 : 20 = 192 : 20 = 9{,}6\text{ (cm)}$.\\
\textbf{Đáp án của câu hỏi là C}
}
\end{ex}

\begin{ex} %[5H5V25]
Một thửa ruộng hình tam giác có độ dài đáy là $25\text{ m}$. Nếu kéo dài đáy thêm $4\text{ m}$ thì diện tích thửa ruộng tăng thêm $36\text{ m}^2$. Diện tích thửa ruộng ban đầu là
\choice
{$200\text{ m}^2$}
{\True $225\text{ m}^2$}
{$250\text{ m}^2$}
{$180\text{ m}^2$}
\loigiai{
Chiều cao của thửa ruộng là:
$36 \times 2 : 4 = 18\text{ (m)}$.\\
Diện tích thửa ruộng ban đầu là:
$25 \times 18 : 2 = 225\text{ (m}^2\text{)}$.\\
\textbf{Đáp án của câu hỏi là B}
}
\end{ex}

\begin{ex} %[5H5V25]
Một hình tam giác có độ dài đáy bằng $18\text{ cm}$ và chiều cao bằng $\dfrac{2}{3}$ độ dài đáy. Diện tích của hình tam giác đó là
\choice
{$216\text{ cm}^2$}
{\True $108\text{ cm}^2$}
{$54\text{ cm}^2$}
{$120\text{ cm}^2$}
\loigiai{
Chiều cao của hình tam giác là:
$18 \times \dfrac{2}{3} = 12\text{ (cm)}$.\\
Diện tích của hình tam giác là:
$18 \times 12 : 2 = 108\text{ (cm}^2\text{)}$.\\
\textbf{Đáp án của câu hỏi là B}
}
\end{ex}

\begin{ex} %[5H5V25]
Nếu tăng độ dài cạnh đáy của một hình tam giác thêm $20\%$ và giữ nguyên chiều cao tương ứng thì diện tích của hình tam giác đó thay đổi như thế nào?
\choice
{Tăng thêm $40\%$}
{Không thay đổi}
{\True Tăng thêm $20\%$}
{Tăng thêm $10\%$}
\loigiai{
Gọi đáy là $a$, chiều cao là $h$. Diện tích ban đầu: $S = \dfrac{a \times h}{2}$.\\
Đáy mới là: $a \times 1{,}2$.\\
Diện tích mới là: $S' = \dfrac{1{,}2 \times a \times h}{2} = 1{,}2 \times S$.\\
Vậy diện tích tăng thêm $20\%$.\\
\textbf{Đáp án của câu hỏi là C}
}
\end{ex}

\begin{ex} %[5H5V25]
Cho tam giác $ABC$ có diện tích bằng $60\text{ cm}^2$. Trên cạnh $BC$ lấy điểm $M$ sao cho $BM = \dfrac{1}{3} BC$. Diện tích tam giác $ABM$ là
\choice
{$15\text{ cm}^2$}
{\True $20\text{ cm}^2$}
{$30\text{ cm}^2$}
{$40\text{ cm}^2$}
\loigiai{
Tam giác $ABM$ và tam giác $ABC$ có chung chiều cao hạ từ đỉnh $A$ xuống đáy $BC$.\\
Vì đáy $BM = \dfrac{1}{3} BC$ nên diện tích tam giác $ABM$ bằng $\dfrac{1}{3}$ diện tích tam giác $ABC$:\\
$S_{ABM} = 60 \times \dfrac{1}{3} = 20\text{ (cm}^2\text{)}$.\\
\textbf{Đáp án của câu hỏi là B}
}
\end{ex}

\begin{ex} %[5H5V25]
Một mảnh đất hình tam giác có diện tích $150\text{ m}^2$. Người ta trồng rau trên mảnh đất đó, cứ mỗi $5\text{ m}^2$ thì thu hoạch được $12\text{ kg}$ rau. Hỏi cả mảnh đất đó thu hoạch được bao nhiêu ki-lô-gam rau?
\choice
{$300\text{ kg}$}
{\True $360\text{ kg}$}
{$400\text{ kg}$}
{$420\text{ kg}$}
\loigiai{
Số ki-lô-gam rau thu hoạch được trên cả mảnh đất là:
$150 : 5 \times 12 = 360\text{ (kg)}$.\\
\textbf{Đáp án của câu hỏi là B}
}
\end{ex}

\begin{ex} %[5H5V25]
Một biển báo giao thông hình tam giác đều có cạnh đáy $60\text{ cm}$ và chiều cao $52\text{ cm}$. Diện tích của biển báo đó là
\choice
{$3\,120\text{ cm}^2$}
{\True $1\,560\text{ cm}^2$}
{$1\,500\text{ cm}^2$}
{$780\text{ cm}^2$}
\loigiai{
Diện tích của biển báo là:
$60 \times 52 : 2 = 1\,560\text{ (cm}^2\text{)}$.\\
\textbf{Đáp án của câu hỏi là B}
}
\end{ex}

\begin{ex} %[5H5V25]
Một hình tam giác có chiều cao $14\text{ cm}$. Nếu tăng chiều cao thêm $3\text{ cm}$ và giữ nguyên độ dài cạnh đáy thì diện tích tam giác tăng thêm $24\text{ cm}^2$. Độ dài cạnh đáy của hình tam giác đó là
\choice
{$12\text{ cm}$}
{\True $16\text{ cm}$}
{$18\text{ cm}$}
{$14\text{ cm}$}
\loigiai{
Phần diện tích tăng thêm là diện tích tam giác có chung đáy với tam giác ban đầu và có chiều cao là $3\text{ cm}$.\\
Độ dài cạnh đáy là:
$24 \times 2 : 3 = 16\text{ (cm)}$.\\
\textbf{Đáp án của câu hỏi là B}
}
\end{ex}

\begin{ex} %[5H5V25]
Cho tam giác $ABC$. Kéo dài cạnh $BC$ về phía $C$ thêm một đoạn $CD = 5\text{ cm}$ thì diện tích tam giác tăng thêm $35\text{ cm}^2$. Biết diện tích tam giác $ABC$ ban đầu là $105\text{ cm}^2$. Độ dài cạnh đáy $BC$ là
\choice
{$12\text{ cm}$}
{\True $15\text{ cm}$}
{$18\text{ cm}$}
{$20\text{ cm}$}
\loigiai{
Chiều cao hạ từ đỉnh $A$ xuống cạnh đáy $BC$ là:
$35 \times 2 : 5 = 14\text{ (cm)}$.\\
Độ dài cạnh đáy $BC$ là:
$105 \times 2 : 14 = 15\text{ (cm)}$.\\
\textbf{Đáp án của câu hỏi là B}
}
\end{ex}

\begin{ex} %[5H5V25]
Cho hình tam giác $ABC$, điểm $M$ là trung điểm của cạnh đáy $BC$. Kết luận nào sau đây là đúng?
\choice
{Diện tích tam giác $ABM$ lớn hơn diện tích tam giác $ACM$}
{Diện tích tam giác $ABM$ bé hơn diện tích tam giác $ACM$}
{\True Diện tích tam giác $ABM$ bằng diện tích tam giác $ACM$}
{Diện tích tam giác $ABC$ gấp ba lần diện tích tam giác $ABM$}
\loigiai{
Hai tam giác $ABM$ và $ACM$ có chung chiều cao hạ từ đỉnh $A$ xuống $BC$ và có đáy $BM = MC$.\\
Do đó, diện tích tam giác $ABM$ bằng diện tích tam giác $ACM$.\\
\textbf{Đáp án của câu hỏi là C}
}
\end{ex}

\begin{ex} %[5H5V25]
Một vườn hoa hình tam giác có độ dài đáy là $15\text{ m}$, chiều cao bằng $\dfrac{4}{5}$ đáy. Người ta dùng $\dfrac{1}{6}$ diện tích mảnh vườn để làm lối đi, phần diện tích còn lại dùng để trồng hoa. Diện tích trồng hoa là
\choice
{$90\text{ m}^2$}
{$15\text{ m}^2$}
{\True $75\text{ m}^2$}
{$60\text{ m}^2$}
\loigiai{
Chiều cao mảnh vườn là:
$15 \times \dfrac{4}{5} = 12\text{ (m)}$.\\
Diện tích mảnh vườn là:
$15 \times 12 : 2 = 90\text{ (m}^2\text{)}$.\\
Diện tích lối đi là:
$90 \times \dfrac{1}{6} = 15\text{ (m}^2\text{)}$.\\
Diện tích trồng hoa là:
$90 - 15 = 75\text{ (m}^2\text{)}$.\\
\textbf{Đáp án của câu hỏi là C}
}
\end{ex}

\textbf{Phần 2. Trắc nghiệm đúng sai}
\setcounterpageref{ex}{0}

\begin{ex} %[5H5V25]
Cho tam giác vuông $ABC$ có hai cạnh góc vuông $AB = 9\text{ cm}, AC = 12\text{ cm}$ và cạnh huyền $BC = 15\text{ cm}$. Kẻ đường cao $AH$ vuông góc với cạnh huyền $BC$.
\choiceTF
{\True Diện tích của tam giác vuông $ABC$ bằng $54\text{ cm}^2$}
{\True Chiều cao $AH$ có độ dài là $7{,}2\text{ cm}$}
{\True Diện tích tam giác $ABC$ bằng một nửa diện tích hình chữ nhật có kích thước $9\text{ cm} \times 12\text{ cm}$}
{Nếu tăng độ dài cạnh $AB$ lên gấp đôi và giữ nguyên cạnh $AC$ thì diện tích tam giác mới gấp $4$ lần diện tích ban đầu}
\loigiai{
\begin{itemchoice}
\item Đúng. Diện tích tam giác: $9 \times 12 : 2 = 54\text{ (cm}^2\text{)}$.
\item Đúng. Chiều cao $AH = 54 \times 2 : 15 = 7{,}2\text{ (cm)}$.
\item Đúng. Vì diện tích tam giác vuông bằng nửa tích hai cạnh góc vuông.
\item Sai. Khi tăng một cạnh góc vuông lên 2 lần thì diện tích chỉ tăng gấp 2 lần.
\end{itemchoice}
}
\end{ex}

\begin{ex} %[5H5V25]
Một thửa ruộng hình tam giác có độ dài cạnh đáy là $30\text{ m}$ và chiều cao bằng $\dfrac{2}{3}$ độ dài đáy.
\choiceTF
{\True Chiều cao của thửa ruộng là $20\text{ m}$}
{\True Diện tích của thửa ruộng là $300\text{ m}^2$}
{Nếu mở rộng đáy thêm $5\text{ m}$ và giữ nguyên chiều cao thì diện tích tăng thêm $100\text{ m}^2$}
{\True Biết rằng cứ $100\text{ m}^2$ thu hoạch được $60\text{ kg}$ lúa thì cả thửa ruộng ban đầu thu hoạch được $180\text{ kg}$ lúa}
\loigiai{
\begin{itemchoice}
\item Đúng. Chiều cao: $30 \times \dfrac{2}{3} = 20\text{ (m)}$.
\item Đúng. Diện tích thửa ruộng: $30 \times 20 : 2 = 300\text{ (m}^2\text{)}$.
\item Sai. Diện tích tăng thêm là: $5 \times 20 : 2 = 50\text{ (m}^2\text{)} 
eq 100\text{ m}^2$.
\item Đúng. Khối lượng lúa thu hoạch được là: $300 : 100 \times 60 = 180\text{ (kg)}$.
\end{itemchoice}
}
\end{ex}

\textbf{Phần 3. Trắc nghiệm trả lời ngắn}
\setcounterpageref{ex}{0}

\begin{ex} %[5H5V25]
Một hình tam giác có diện tích bằng $72\text{ cm}^2$ và chiều cao tương ứng là $9\text{ cm}$. Độ dài cạnh đáy của hình tam giác đó là bao nhiêu xăng-ti-mét?
\shortans{16}
\loigiai{
Độ dài cạnh đáy của hình tam giác là:
$72 \times 2 : 9 = 16\text{ (cm)}$.
}
\end{ex}

\begin{ex} %[5H5V25]
Một mảnh đất hình tam giác có diện tích $120\text{ m}^2$. Nếu kéo dài cạnh đáy thêm $3\text{ m}$ thì diện tích mảnh đất tăng thêm $18\text{ m}^2$. Độ dài cạnh đáy ban đầu của mảnh đất là bao nhiêu mét?
\shortans{20}
\loigiai{
Chiều cao của mảnh đất là:
$18 \times 2 : 3 = 12\text{ (m)}$.\\
Độ dài cạnh đáy ban đầu là:
$120 \times 2 : 12 = 20\text{ (m)}$.
}
\end{ex}

\begin{ex} %[5H5V25]
Cho tam giác $ABC$ có diện tích bằng $96\text{ cm}^2$. Trên cạnh $BC$ lấy điểm $M$ sao cho $BM = 3 \times MC$. Tính diện tích tam giác $ABM$ theo đơn vị xăng-ti-mét vuông.
\shortans{72}
\loigiai{
Vì $BM = 3 \times MC$ nên đoạn $BC$ gồm $4$ phần bằng nhau và $BM$ chiếm $3$ phần, tức $BM = \dfrac{3}{4} BC$.\\
Hai tam giác $ABM$ và $ABC$ có chung đường cao kẻ từ đỉnh $A$ nên:
$$S_{ABM} = \dfrac{3}{4} \times S_{ABC} = \dfrac{3}{4} \times 96 = 72\text{ (cm}^2\text{)}.$$
}
\end{ex}

\begin{ex} %[5H5V25]
Một tấm tôn hình tam giác có cạnh đáy dài $2{,}4\text{ m}$ và chiều cao bằng $1{,}5\text{ m}$. Người ta cắt bỏ một phần tôn có diện tích bằng $0{,}3\text{ m}^2$. Diện tích phần tôn còn lại là bao nhiêu đề-xi-mét vuông?
\shortans{150}
\loigiai{
Diện tích tấm tôn ban đầu là:
$2{,}4 \times 1{,}5 : 2 = 1{,}8\text{ (m}^2\text{)}$.\\
Diện tích phần tôn còn lại là:
$1{,}8 - 0{,}3 = 1{,}5\text{ (m}^2\text{)}$.\\
Đổi: $1{,}5\text{ m}^2 = 150\text{ dm}^2$.
}
\end{ex}

\begin{ex} %[5H5V25]
Một hình tam giác có độ dài cạnh đáy bằng $25\text{ cm}$. Nếu tăng độ dài cạnh đáy thêm $5\text{ cm}$ thì diện tích hình tam giác tăng thêm $40\text{ cm}^2$. Tính diện tích ban đầu của hình tam giác đó theo đơn vị xăng-ti-mét vuông.
\shortans{200}
\loigiai{
Chiều cao của tam giác là:
$40 \times 2 : 5 = 16\text{ (cm)}$.\\
Diện tích ban đầu của hình tam giác là:
$25 \times 16 : 2 = 200\text{ (cm}^2\text{)}$.
}
\end{ex}

\begin{ex} %[5H5V25]
Cho hình chữ nhật $ABCD$ có chiều dài $AB = 24\text{ cm}$ và chiều rộng $BC = 15\text{ cm}$. Lấy một điểm $M$ bất kì nằm trên cạnh $CD$. Diện tích hình tam giác $MAB$ là bao nhiêu xăng-ti-mét vuông?
\shortans{180}
\loigiai{
Hình tam giác $MAB$ có đáy $AB = 24\text{ cm}$ và chiều cao hạ từ đỉnh $M$ xuống đáy $AB$ chính bằng chiều rộng $BC = 15\text{ cm}$.\\
Diện tích tam giác $MAB$ là:
$24 \times 15 : 2 = 180\text{ (cm}^2\text{)}$.
}
\end{ex}
